% 第二节：张量与指标缩并
\chapter{张量与指标缩并}

\section{指标的上下与自由指标}

第一章说明了：只有矢量与对偶矢量之间才有自然的缩并 $f\cdot A=f_iA^i$，这一运算与坐标系的
选取无关。后文要用一个二阶张量表示圆锥曲线，为此先把矢量与对偶矢量的分量记号推广到
一般的张量。

张量不依赖于坐标系的定义要借助多重线性映射，比较抽象。这里始终假定已取定一组基底，
用分量的观点来处理张量；这样得到的记号与坐标无关的说法是一致的，而在运算中更为直接。

设 $\{e_i\}$ 是矢量空间的一组基，矢量 $A$ 在该基底下写成
\[
A=\sum_i A^i e_i ,
\]
$A^i$ 称为 $A$ 的分量；设 $\{e^i\}$ 是对偶基，对偶矢量 $f$ 写成 $f=\sum_i f_i e^i$，$f_i$
称为 $f$ 的分量。为了在记号上区分二者，作如下约定：

\begin{definition}[上下指标的约定]
分量指标满足下列约定：
\begin{itemize}
\item 矢量的分量指标写在右上，称为上指标；
\item 对偶矢量的分量指标写在右下，称为下指标；
\item 单独出现的指标默认取遍 $1,\dots,n$，称为自由指标；
\item 给出 $n$ 个实数 $A^i$（或 $f_i$），就认为给出了相应的矢量（或对偶矢量）。
\end{itemize}
\end{definition}

在上下指标的记号下，$f$ 作用于 $A$ 得到
\[
f\cdot A=f_1A^1+\cdots+f_nA^n=\sum_i f_iA^i .
\]
也就是说：上指标与下指标之间可以缩并求和，而两个上指标之间或两个下指标之间不能这样求和。

\section{爱因斯坦求和约定}

在书写中，求和号常常不便。因此约定：

\begin{definition}[爱因斯坦求和约定]
在一个表达式中，若某个指标字母恰好出现两次，并且一次在上、一次在下，则称它为哑指标，
并约定对该指标从 $1$ 到 $n$ 求和，求和号省略不写。只出现一次的指标称为自由指标。
\end{definition}

既然哑指标意味着求和，用哪个字母代表哑指标就没有区别：$f_iA^i$ 与 $f_kA^k$ 表示同一个数。
例如 $f_iA^jB^i$ 表示一个矢量，其中 $i$ 是哑指标、$j$ 是自由指标，它的第 $j$ 个分量是
\[
(f_1B^1+\cdots+f_nB^n)A^j .
\]

需要注意的是，形如 $A^iB^i$、$A^iB_iC^iD_i$、$A^iB_iC_i$ 的记号没有意义：前两者中同一个
指标出现了两次以上，后者中 $i$ 两次都写在下方。因此，像
$A^1B_1C_1+\cdots+A^nB_nC_n$ 这种逐分量相乘再求和的式子不能直接写成求和约定的形式；
不过由几何的坐标无关性，后面需要的运算绝大多数都能用求和约定表示。

\section{张量的分量与类型}

$f_iA^j$ 有两个自由指标 $i,j$，各自可取 $1,\dots,n$，因此它由 $n^2$ 个实数组成。有 $m$ 个
自由指标的对象称为 $m$ 阶张量；二阶张量一般写成 $M_i{}^j$，含义是每一组 $(i,j)$ 对应一个
实数作为分量，其中 $i$ 是下指标、$j$ 是上指标。

\begin{definition}[(k,l) 型张量]
有 $k$ 个上指标、$l$ 个下指标的 $k+l$ 阶张量称为 $(k,l)$ 型张量。
\end{definition}

例如 $A^{ij}{}_k{}^l$ 是 $(3,1)$ 型张量。一个张量的型不仅取决于上、下指标的个数，指标排列
的顺序也重要；代表指标的字母则无关紧要。

区分上、下指标是必要的。例如设 $M_i{}^j=f_iA^j$，则 $M_i{}^j$ 的上指标可以与对偶矢量
$u_j$ 缩并：
\[
M_i{}^ju_j=f_iA^ju_j ,
\]
结果是一个对偶矢量；而形如 $M_{ij}$ 的张量不能与 $u_j$ 缩并，$f_iA_ju_j$ 是没有意义的。
一般地，一个指标是上指标还是下指标，决定了它能与对偶矢量还是矢量缩并。

二阶张量容易与矩阵混淆，因为二者都有两个指标。区别在于：张量是几何概念，矩阵是代数概念。
借助坐标系可以把二阶张量写成矩阵（正如把矢量写成有序数组），但意义不同：矩阵不区分上、
下指标，把不同类型的二阶张量写成矩阵会丢失类型信息。

\section{缩并}

\begin{definition}[缩并]
设 $T$ 是 $(k,l)$ 型张量，$k\ge1$，$l\ge1$。取出它的一个上指标与一个下指标，令二者取同一个
字母并对该指标求和，所得结果是 $(k-1,l-1)$ 型张量。这一运算称为缩并。
\end{definition}

例如对 $A^{ij}{}_k{}^l$ 的第 $1$ 个上指标 $i$ 与第 $1$ 个下指标 $k$ 缩并，得到
\[
B^{jl}=A^{ij}{}_i{}^l ,
\]
其中 $j,l$ 仍是自由指标。缩并是矢量与对偶矢量之间缩并的推广。

\section{(1,1) 型张量与矩阵}

$(1,1)$ 型张量的性质可以完全用矩阵的语言复述。

\begin{proposition}[(1,1) 型张量的五条性质]
设 $M^i{}_j$ 是 $(1,1)$ 型张量，则
\begin{enumerate}
\item 它把矢量映为矢量：$B^i=M^i{}_jA^j$，因此是矢量空间到自身的线性映射；
\item 它把对偶矢量映为对偶矢量：$B_j=M^i{}_jA_i$；
\item 两个 $(1,1)$ 型张量可以自然地相乘，乘积仍是 $(1,1)$ 型张量：
$M^i{}_k=A^i{}_jB^j{}_k$；
\item 存在 $(1,1)$ 型张量 $\delta^i_j$，它作用于矢量得到该矢量本身，作用于对偶矢量也得到该
对偶矢量本身，与任意 $(1,1)$ 型张量相乘仍得到该张量，称为恒等张量；
\item 它的上指标与下指标可以缩并，得到一个实数 $M^i{}_i$，称为 $M$ 的迹。
\end{enumerate}
\end{proposition}

这五条性质依次对应于矩阵与列矢量相乘、行矢量与矩阵相乘、矩阵乘法、单位矩阵与矩阵的迹。
因此，把矩阵看作线性变换时，矩阵乘法正对应于 $(1,1)$ 型二阶张量的运算。

\begin{remark}
把第 (1)(2) 条分别称为矩阵的“左乘”与“右乘”，依赖于把 $M^i{}_j$ 写成矩阵时约定 $i$ 为行
指标、$j$ 为列指标：此时 $B^i=M^i{}_jA^j$ 就是矩阵乘列矢量，而 $B_j=M^i{}_jA_i$ 就是行矢量
乘矩阵。
\end{remark}

在线性代数中，矩阵还可以用来表示二次型。二次型那样的对象不是 $(1,1)$ 型张量，而更接近
$(0,2)$ 型或 $(2,0)$ 型张量，这是下一章的内容。

\begin{exercise}
把 $f_iA^iB^j$ 按求和约定展开，指出其中的哑指标与自由指标，并说明结果是什么类型的量，
它的分量是什么。
\end{exercise}
